\part{Lösungsansätze}
\begin{frame}
    \partpage
\end{frame}
\section{Kantenansatz}
\begin{frame}{Kantenansatz}
\begin{figure}[h]
    \centering
    \includegraphics[width=0.7\linewidth]{beamer/Kanten_erklärung.pdf}
\end{figure}  
\end{frame}

\note[itemize]{
    \item vollständiger Graph erstellt
    \item Kanten kein Knoten schneiden
    \item jede Kante Variable
    \item _
    \item notwendige Bedingungen
    \item optionale Bedingungen
    \item 
}

\begin{frame}{Überschneidungsbedingung}
\begin{figure}[h]
    \centering
    \includegraphics[width=0.6\linewidth]{beamer/Überschneidung.pdf}
\end{figure}  
\end{frame}

\note[itemize]{
    \item Kanten dürfen sich nicht schneiden
    \item Kreuzende Paare finden
    \item nur eine Kante pro Paar
}

\begin{frame}{Gradbeschränkung}
\begin{figure}[h]
    \centering
    \includegraphics[width=0.6\linewidth]{beamer/Gradbeschränkung.pdf}
\end{figure}  
\end{frame}

\note[itemize]{
    \item inzidente Kanten gleich Sollgrad
    \item Sollgrad 3
    \item 3 von 6 aktiv
}

\begin{frame}{Maximale Kanten}
    \fabian{Formal raus?}
    % \pageText{\(6 n - 6 - 2k \overset{!}{=} \sum \degrees(i)\)}
    \pageText{\(\text{Anzahl Kanten} \cdot 2 \overset{!}{=} \sum \text{Sollgrad}\)}
\end{frame}

\note[itemize]{
    \item Validieren das Sollgrade möglich
    \item Anzahl Kanten: \(3n - 3 - k\)
    \begin{itemize}
        \item n: alle Knoten
        \item k: Knoten auf Hülle
    \end{itemize}
    \item wenn nicht stimmt keine Lösung möglich
    \item testet indirekt Kanten maximal
}

\begin{frame}{Gradbeschränkung}
    \fabian{Durchschnitt}
\end{frame}

\note[itemize]{
    \item Zielfunktion kein Gradbeschränkung
    \item Maximal zusätzlich Prüfen
    \item _
    \item Differenz zwischen Soll und Istgrad
    \item Differenz durch Sollgrad für Prozentuale abweichung
    \item Summe durch Anzahl
    \item Durchschnitt der Prozentualen Abbildung
}


\begin{frame}{Gradbeschränkung}
    \fabian{Max min}
\end{frame}

\note[itemize]{
    \item Maximale Abweichung Minimieren
    \item var mindestens groß wie jede abweichung
    \item diese var minimieren
}

\begin{frame}{alternative Kantenbedingung}
    \fabian{sollte ich hier eine Refernz auf euer Paper machen?}
    \begin{figure}[h]
    \centering
    \includegraphics[width=0.7\linewidth]{all_edges.pdf}
\end{figure}  
\end{frame}

\note[itemize]{
    \item maximiert Kanten
    \item für jede Kante gild
    \begin{itemize}
        \item Kante ist aktiv
        \item oder Schnittkante aktive
    \end{itemize}
    \item für gegebenen Graphen alle Triangulierungen
    \item immer e1 aktiv oder schnittkante
}

\begin{frame}{Hülle Fixieren}
    \begin{figure}[h]
    \centering
    \includegraphics[width=0.6\linewidth]{beamer/Hülle.pdf}
\end{figure}  
\end{frame}



\begin{frame}{Kanten Fixieren / Ausschließen}
    \begin{figure}[h]
    \centering
    \includegraphics[width=0.8\linewidth]{exclude_edges/ecludegrade2.pdf}
\end{figure}  
\end{frame}

\begin{frame}{Kanten Ausschließen}
    \begin{minipage}{0.48\textwidth}
        \centering
        \includegraphics[width=0.8\textwidth]{exclude_edges/exclude_edge_0.pdf}
    \end{minipage}
    \hfill
    \begin{minipage}{0.48\textwidth}
        \centering
        \includegraphics[width=0.8\textwidth]{exclude_edges/exclude_edge_1.pdf}
    \end{minipage}
\end{frame}

\begin{frame}{Kanten Ausschließen}
    \begin{figure}[h]
    \centering
    \includegraphics[width=0.7\linewidth]{exclude_edges/exclude_edge_intersection.pdf}
\end{figure}  
\end{frame}

\section{Dreiecksansatz}
\begin{frame}{Dreiecksansatz}
\begin{figure}[h]
    \centering
    \includegraphics[width=0.6\linewidth]{beamer/Dreieck_erklärung.pdf}
\end{figure}  
\end{frame}

\begin{frame}{Gradbeschränkung}
    \begin{minipage}{0.56\textwidth}
        \centering
        \includegraphics[width=0.8\textwidth]{beamer/Gradbeschränkung_tri.pdf}
    \end{minipage}
    \hfill
    \begin{minipage}{0.40\textwidth}
        \centering
        \includegraphics[width=0.8\textwidth]{beamer/Gradbeschränkung_tri_edge.pdf}
    \end{minipage}
\end{frame}

\begin{frame}{Überschneidungsbedingung}
\begin{figure}[h]
    \centering
    \includegraphics[width=0.7\linewidth]{beamer/Überschneidung_tri.pdf}
\end{figure} 
\end{frame}

\begin{frame}{Überschneidungsbedingung}
\begin{figure}[h]
    \centering
    \includegraphics[width=0.6\linewidth]{beamer/Überschneidung_tri_edge.pdf}
\end{figure} 
\end{frame}

\begin{frame}{Dreiecks Ausschließen}
    \begin{figure}[h]
    \centering
    \includegraphics[width=0.8\linewidth]{beamer/Dreieck_ausschließen.pdf}
\end{figure} 
\end{frame}

\begin{frame}{Subset}
    \fabian{ich habe jetzt subset in den Diagrammen drin, ist aber schlecht. Soll ich das dann hier erklären oder wie mache ich das am besten? Einfach nur auf die Arbeit verweisen?}
\end{frame}

